\bibitem{buneman2001} P. Buneman, S. Khanna, and W.-C. Tan. Why and Where: A Characterization of Data Provenance. In \emph{Database Theory---ICDT 2001}, pages 316--330. Springer, 2001. doi:10.1007/3-540-44503-X_20.

\bibitem{green2007} T. J. Green, G. Karvounarakis, and V. Tannen. Provenance semirings. In \emph{Proceedings of PODS 2007}, pages 31--40, 2007. doi:10.1145/1265530.1265535.

\bibitem{prov2013} W3C. \emph{PROV-O: The PROV Ontology}. W3C Recommendation, 30 April 2013.

\bibitem{wilkinson2016} M. D. Wilkinson et al. The FAIR Guiding Principles for scientific data management and stewardship. \emph{Scientific Data} 3:160018, 2016. doi:10.1038/sdata.2016.18.

\bibitem{zhang2020claim} Y. Zhang, Z. Ives, and D. Roth. ``Who said it, and Why?'' Provenance for Natural Language Claims. In \emph{Proceedings of ACL 2020}, pages 4416--4426. doi:10.18653/v1/2020.acl-main.406.

\bibitem{boelling2014see} C. B\"olling, M. Weidlich, and H.-G. Holzh\"utter. SEE: structured representation of scientific evidence in the biomedical domain using Semantic Web techniques. \emph{Journal of Biomedical Semantics} 5:S1, 2014. doi:10.1186/2041-1480-5-S1-S1.

\bibitem{amaral2024prove} G. Amaral, O. Rodrigues, and E. Simperl. ProVe: A pipeline for automated provenance verification of knowledge graphs against textual sources. \emph{Semantic Web} 15(6):2159--2192, 2024. doi:10.3233/SW-233467.

\bibitem{vladika2024} J. Vladika and F. Matthes. Comparing Knowledge Sources for Open-Domain Scientific Claim Verification. In \emph{Proceedings of EACL 2024}, pages 2103--2114. doi:10.18653/v1/2024.eacl-long.128.

\bibitem{martinboyle2026} A. Martin-Boyle, C. A. C. Leckey, M. C. Brown, and H. Kaur. PaperTrail: A Claim-Evidence Interface for Grounding Provenance in LLM-based Scholarly Q\&A. arXiv:2602.21045, 2026.

\bibitem{ramos2021eventclaim} J. Ramos, J. Rasga, C. Sernadas, and L. Vigan\`o. Event-Based Time-Stamped Claim Logic. \emph{Journal of Logical and Algebraic Methods in Programming} 121:100684, 2021. doi:10.1016/j.jlamp.2021.100684.

\bibitem{donto2026} T. Davis and A. Davis. donto: An Evidence Operating System for Contested Knowledge. Systems report, 28 May 2026.

\bibitem{vitali2026dec} F. Vitali and V. Pasqual. Provenance-Enhanced Statements in Knowledge Graphs. arXiv:2606.15246, 2026.

\bibitem{mccarthy1993} J. McCarthy. Notes on Formalizing Context. In \emph{Proceedings of the 13th International Joint Conference on Artificial Intelligence}, 1993.

\bibitem{giunchiglia1993} F. Giunchiglia. Contextual reasoning. \emph{Epistemologia}, 1993.

\bibitem{abramsky2011} S. Abramsky and A. Brandenburger. The sheaf-theoretic structure of non-locality and contextuality. \emph{New Journal of Physics} 13:113036, 2011. doi:10.1088/1367-2630/13/11/113036.

\bibitem{curry2014} J. Curry. \emph{Sheaves, Cosheaves and Applications}. PhD thesis, University of Pennsylvania, 2014. arXiv:1303.3255.

\bibitem{robinson2017} M. Robinson. Sheaves are the canonical data structure for sensor integration. \emph{Information Fusion} 36:208--224, 2017. doi:10.1016/j.inffus.2016.12.002.

\bibitem{dechter1992} R. Dechter. From local to global consistency. \emph{Artificial Intelligence} 55(1):87--107, 1992. doi:10.1016/0004-3702(92)90043-W.

\bibitem{pearl2009} J. Pearl. \emph{Causality: Models, Reasoning, and Inference}. 2nd edition, Cambridge University Press, 2009.

\bibitem{halpern2016} J. Y. Halpern. \emph{Actual Causality}. MIT Press, 2016.

\bibitem{birkhoff1940} G. Birkhoff. \emph{Lattice Theory}. American Mathematical Society Colloquium Publications, 1940.

\bibitem{oles2000} F. J. Oles. An application of lattice theory to knowledge representation. \emph{Theoretical Computer Science} 249(1):163--196, 2000. doi:10.1016/S0304-3975(00)00058-X.

\bibitem{ganter1999} B. Ganter and R. Wille. \emph{Formal Concept Analysis: Mathematical Foundations}. Springer, 1999.

\bibitem{yanglee2026fca} Y. Yang and H. Lee. Verifiable Knowledge Expansion through Retrieval-Grounded Formal Concept Analysis. arXiv:2607.01773, 2026.

\bibitem{cousot1977} P. Cousot and R. Cousot. Abstract Interpretation: A Unified Lattice Model for Static Analysis of Programs by Construction or Approximation of Fixpoints. In \emph{Proceedings of POPL 1977}, pages 238--252, 1977. doi:10.1145/512950.512973.

\bibitem{tarski1955} A. Tarski. A lattice-theoretical fixpoint theorem and its applications. \emph{Pacific Journal of Mathematics} 5(2):285--309, 1955.
